\documentclass[aspectratio=169]{ctexbeamer}
\usepackage[T1]{fontenc}
\usepackage{latexsym,amsmath,xcolor,multicol,booktabs,calligra}
\usepackage{graphicx,listings,stackengine}
\usepackage{hyperref}

\usepackage{PekingU}

% 去掉每节/子节前面的自动目录页
\AtBeginSection[]{}
\AtBeginSubsection[]{}

% ===== 每节课改这里 =====
\author{Cap1taL}
\title{树上问题 1}
\subtitle{我听说，从前有个人，闲来无事，想找个人一起下棋。他站在院子里，喊了一声：“谁有空与我手谈一局？”结果邻居们各有各的事，有的要浇园，有的要喂鸡，有的在补衣裳。他便垂着手站在门口，等了一会儿，又喊了一遍。这时有个路过的人说：“你喊了也没用，大家都有自己的本分。”那人听了，想了想，觉得也对，便自己回屋摆开棋盘，左手与右手对弈，倒也下出了一番道理。如今我喊一声“谁有空跟我打游戏”，其实也是一样的道理。游戏这东西，本是人与人之间的交际，讲的是你来我往的礼数。若无人应答，我便该明白，此刻大家各有各的职分，各有各的忙处。我不该强求，也不必失落。自己开一局，练练手，等有人得空了，自然有人应声。这就可以说是：有伴便同乐，无伴便独行，二者都不失去体面。}
\date{2026年8月}
% ========================

% 代码高亮设置
\definecolor{deepblue}{rgb}{0,0,0.5}
\definecolor{deepred}{rgb}{0.6,0,0}
\definecolor{deepgreen}{rgb}{0,0.5,0}
\definecolor{halfgray}{gray}{0.55}

\lstset{
    basicstyle=\ttfamily\small,
    keywordstyle=\bfseries\color{deepblue},
    emphstyle=\ttfamily\color{deepred},
    stringstyle=\color{deepgreen},
    numbers=left,
    numberstyle=\small\color{halfgray},
    rulesepcolor=\color{red!20!green!20!blue!20},
    frame=shadowbox,
}

\begin{document}

\kaishu

% 封面
\begin{frame}
    \titlepage
\end{frame}

% 目录
% \begin{frame}{目录}
%     \tableofcontents[sectionstyle=show,subsectionstyle=show/shaded/hide,subsubsectionstyle=show/shaded/hide]
% \end{frame}

% ===== 从这里开始写你的内容 =====

\section{基本算法}

\subsection{树上背包DP}
\begin{frame}{树上背包 DP 的复杂度}
    \begin{itemize}
        \item 树上背包转移时枚举子树大小，需要记录前缀子树的子树和，而非真正的 $sz_u$
        \item 复杂度证明考虑所有的点对
        \item 在容量与 $n$ 同阶时为 $O(n^2)$，在 $m<<n$ 时为 $O(nm)$
    \end{itemize}
\end{frame}

\subsection{树链剖分}
\begin{frame}{树链剖分}
    \begin{itemize}
        \item 树链剖分可以把一棵树划分成若干链（自上而下的链），以便于维护树上信息
        \item 树链剖分一般会对链内重标号，并保证一条链内点的编号连续。
        \item 因此，我们可以用维护序列上的来维护一条链上的信息
        \item 对于跨链的信息（例如路径），我们一般将其拆分成若干链。合适的树链剖分方式往往可以保证链的数量较少，从而保持较优的复杂度
    \end{itemize}
\end{frame}

\begin{frame}{重链剖分}
    \begin{itemize}
        \item Code \url{https://www.luogu.com.cn/paste/vqjegxoq}
        \item 选择最重的儿子进行链剖分延伸
        \item 链切换次数为 $O(\log n)$，证明考虑每次换链一定子树大小翻倍
        \item 链上编号连续，子树编号也连续
        \item 可以把一个路径拆为 $O(\log n)$ 个序列上的区间，或者维护子树信息
        \item 常数很小
    \end{itemize}
\end{frame}

\begin{frame}{长链剖分}
    \begin{itemize}
        \item Code \url{https://www.luogu.com.cn/paste/vqjegxoq}
        \item 选择最深的儿子进行链剖分延伸
        \item 链切换次数为 $O(\sqrt{n})$，证明考虑每次换链一定多一条更长的链
        \item 由于链切换次数较劣，我们一般不用长链剖分维护路径或子树等树上信息
    \end{itemize}
\end{frame}

\begin{frame}{长链剖分优化 DP}
    \begin{itemize}
        \item 类似树上启发式合并的一种 trick
        \item 可以把与深度有关的信息维护加速到线性
        \item 暴力做这类信息维护一般是 $O(n^2)$ 的
        \item 具体来说，每次直接继承长儿子的信息，暴力合并轻儿子的信息
        \item 继承是快速的，一个链只有在链顶向上合并的时候才需要完整遍历，总复杂度 $O(n)$
        \item 由于不能完整开数组，实现一般采用指针+动态分配内存，或者使用 vector 与 swap。
    \end{itemize}
\end{frame}

\begin{frame}{长链剖分求 k 级祖先}
    \begin{itemize}
        \item 注意到关键性质：一个点 $k$ 级祖先所在的长链长度必 $ge k$
        \item 反证法，若小于 $k$，则选取当前这个链更优，这个才是长链
        \item 因此若我们先跳一次 $2^h$ 祖先，剩余高度不会超出所在链的预处理范围
        \item 对于每个链顶，设链长为 $L$，处理其上下各 $L$ 个点，跳一次 $2^h$ 后直接处理即可
        \item 预处理 $O(n\log n)$，查询 $O(1)$
        \item 有时仍然不如重链剖分快
    \end{itemize}
\end{frame}

\subsection{最近公共祖先 LCA}
\begin{frame}{树上倍增求 LCA}
    \begin{itemize}
        \item 预处理一个点的 $2^k$ 级祖先
        \item 求 LCA 的时候先跳平，再从大到小尝试跳到 LCA 下方
        \item 预处理 $O(n\log n)$，查询 $O(\log n)$，空间 $O(n\log n)$，实际运行超级慢
    \end{itemize}
\end{frame}

\begin{frame}{重链剖分求 LCA}
    \begin{itemize}
        \item 预处理重链剖分，不需要重标号，只需要处理链顶
        \item 不断跳重链，直到跳到同一个链上，较浅的一个即为 LCA
        \item 预处理 $O(n)$，查询 $O(\log n)$，空间 $O(n)$，实际运行超级快
    \end{itemize}
\end{frame}

\begin{frame}{DFN 序求 LCA}
    \begin{itemize}
        \item 考虑 $u,v$ 不是祖先后代的情况。显然 dfn 序列上 $u,v$ 之间的元素都在 $lca$ 的子树中。只需要找到深度最小的元素，其父亲即为 LCA。这是一个 RMQ 问题，可以使用 ST 表。
        \item 对于他们是祖先后代的情况，只需要去掉 $u$，即可兼容上一种情况。
        \item 具体实现并不需要维护父亲与深度，只需要在 ST 表底层直接填写父亲，找 dfn 序最小的即可。容易理解这个技巧的正确性
        \item Code：\url{https://www.luogu.com.cn/record/152919187}
        \item 预处理 $O(n\log n)$，查询 $O(1)$，空间 $O(n\log n)$，常数远小于树上倍增欧拉环游序
    \end{itemize}
\end{frame}

\subsection{树分治}
\begin{frame}{点分治}
    \begin{itemize}
        \item 又称重心分解，可以用来处理大规模的路径问题等
        \item 我们在树上选择一个根，则路径可以按照是否经过与经过的情况，被分为三种
        \item 若我们能适当地选取根，则可以高效地批量分类并处理这些路径
        \item 每次树根，暴力遍历整个极大联通块内的各个子树，维护信息的同时求出后续的新重心，复杂度是 $O(n\log n)$ 的
    \end{itemize}
\end{frame}
\begin{frame}{点分树}
    \begin{itemize}
        \item 把点分治的形态固定成一棵树，属于一种重构树
        \item 可以处理一类与树形态无关的问题
        \item 关键问题在于，如何去除点分树上子树对父节点的影响
        \item 常见做法是对于每个节点维护两个数据结构，代表字数内点与自己和与父亲的信息
        \item 另一种做法是维护前缀与后缀数据结构，常见于修改极其特殊的情形
    \end{itemize}
\end{frame}




\section{习题}

\subsection{P4315 月下“毛景树”}
\begin{frame}[fragile]{P4315 月下“毛景树”}
    给定一棵有 $N$ 个节点、$N-1$ 条树枝的树，毛毛果长在树枝上。支持以下操作：
    \begin{verbatim}
Change k w   : 将第 k 条树枝的毛毛果个数改为 w
Cover u v w  : 将 u 到 v 路径上所有树枝的毛毛果个数改为 w
Add u v w    : 将 u 到 v 路径上所有树枝的毛毛果个数增加 w
Max u v      : 查询 u 到 v 路径上树枝毛毛果个数的最大值
    \end{verbatim}
    操作以 \texttt{Stop} 结束。

    \medskip
    $1\le N\le 10^5$，操作和询问数目不超过 $10^5$。
\end{frame}

\begin{frame}{Solution}
    \begin{itemize}
        \item 树剖模板
        \item Code：\url{https://www.luogu.com.cn/record/109990824}
    \end{itemize}
\end{frame}
\subsection{CF1009F Dominant Indices}

\begin{frame}{CF1009F Dominant Indices}
    给定一棵以 $1$ 为根的树，包含 $n$ 个顶点。
    对于每个顶点 $x$，定义 $d_{x,i}$ 为：在 $x$ 的子树中，与 $x$ 距离恰好为 $i$ 的顶点数。
    顶点 $x$ 的主导下标 $j$ 是满足以下条件的最小下标：
    \[
        d_{x,j}=\max_{k\ge 0} d_{x,k}.
    \]
    换句话说，$j$ 是 $d_{x,j}$ 第一次取到最大值的深度。
    请计算每个顶点的主导下标。

    \medskip
    $1\le n\le 10^6$。
\end{frame}

\begin{frame}{Solution}
    \begin{itemize}
        \item 可以写出一个暴力 DP。设 $f_{u,i}$ 表示 $u$ 这个点下面深度为 $i$ 的点数
        $$
            f_{u,i}={\sum_{u\to v}}f_{v,i-1}
        $$
        \item 暴力做显然是 $n^2$ 的
        \item 采用长链剖分优化。对于每个点：
        \begin{itemize}
            \item 直接继承长儿子的 DP 数组
            \item 在数组最前面添加一个 $ 1 $ 代表 $f_{u,0}$
            \item 暴力合并所有短儿子
            \item 复杂度 $O(n)$
        \end{itemize}
    \end{itemize}
\end{frame}

\subsection{P3806 【模板】点分治}

\begin{frame}{P3806 【模板】点分治}
    给定一棵有 $n$ 个点的树，边有边权。
    共有 $m$ 次询问，每次询问一个整数 $k$，判断树上是否存在一对点，使得它们之间的距离恰好为 $k$。

    \medskip
    $1\le n\le 10^4$，$1\le m\le 100$，$1\le k\le 10^7$，$1\le w\le 10^4$。
\end{frame}

\begin{frame}{Solution}
    \begin{itemize}
        \item 点分治
        \item 每次分治时遍历子树，记录到根的距离，等价于拆点
        \item 使用某些数据结构或者桶支持查询即可
        \item 某种意义上类似树上 DP 的转移
        \item 询问数较少，可以暴力处理。复杂度 $O(nm\log n)$
        \item Code：\url{https://www.luogu.com.cn/record/291033843}
    \end{itemize}
\end{frame}

\subsection{P6329 【模板】点分树 / 震波}
\begin{frame}[fragile]{P6329 【模板】点分树 / 震波}
    给定一棵 $n$ 个点的树，相邻两个城市距离为 $1$，第 $i$ 个城市价值为 $value_i$。
    在线处理 $m$ 次操作：
    \begin{verbatim}
0 x k：查询与 x 距离不超过 k 的所有城市的价值和。
1 x y：将第 x 个城市的价值改为 y。
    \end{verbatim}
    强制在线

    \medskip
    $1\le n,m\le 10^5$，$1\le value_i,y\le 10^4$，$0\le k\le n-1$。
\end{frame}

\begin{frame}{Solution}
    \begin{itemize}
        \item 建立点分树
        \item 每个点维护两个数据结构，存储 点分树子树 内与自己 \ 点分树父亲距离 $\le k$ 的价值信息
        \item 查询时暴力在点分树上上跳，每次计算都扣除子树影响
        \item 修改同样暴力跳点分树，维护两个数据结构
        \item 此题的实现可以使用树状数组，时间复杂度 $O(n\log^2n)$，空间复杂度 $O(n\log n)$
    \end{itemize}
\end{frame}

\subsection{P2056 [ZJOI2007] 捉迷藏}
\begin{frame}[fragile]{P2056 [ZJOI2007] 捉迷藏}
    给定一棵由 \(N\) 个房间和 \(N-1\) 条双向走廊组成的树，初始所有房间的灯都是关闭的。
    孩子们只会躲藏在没有开灯的房间中。每次操作为：
    \begin{verbatim}
C i : 改变第 i 个房间的开关状态（开 <-> 关）
G   : 查询最远的两个关灯房间的距离
    \end{verbatim}
    对于每个 \texttt{G} 操作：若只有一间关灯房间，输出 \(0\)；若没有关灯房间，输出 \(-1\)。

    \medskip
    \(1\le N\le 100000\)，\(1\le Q\le 500000\)。
\end{frame}

\begin{frame}{Solution}
    \begin{itemize}
        \item 建立点分树
        \item 每个点维护两个数据结构，存储 点分树子树 内与自己 \ 点分树父亲的距离信息
        \item 查询时暴力在点分树上上跳，每次计算都扣除子树影响
        \item 修改同样暴力跳点分树
        \item 具体实现可以维护两类可删堆，一类是存储所有路径的堆，一类是存储各个子树最大值的堆。或者使用线段树等。总之要维护子树到根路径信息
        \item 时间复杂度 $O(n\log^2n)$，空间复杂度 $O(n\log n)$
    \end{itemize}
\end{frame}

\subsection{P9058 [Ynoi2004] rpmtdq}
\begin{frame}[fragile]{P9058 [Ynoi2004] rpmtdq}
    给定一棵有边权的无根树，共有 \(n\) 个节点。
    定义 \(\mathrm{dist}(i,j)\) 表示树上点 \(i\) 和点 \(j\) 之间的距离。

    有 \(q\) 次询问，每次给出 \(l,r\)，需要求出：
    \[
    \min_{l\le i<j\le r} \mathrm{dist}(i,j)
    \]
    若不存在满足 \(l\le i<j\le r\) 的二元组 \((i,j)\)，则输出 \(-1\)。

    \medskip
    \(1\le n\le 2\times 10^5\)，\(1\le q\le 10^6\)，\(1\le z\le 10^9\)。
\end{frame}

\begin{frame}{Solution}
    \begin{itemize}
        \item 寻找支配点对
        \item 固定分治根，考虑什么样的点对无效，即对于 $i<j<k$，若 $dis_{rt,k}>dis{rt_j}$，则 $k$ 不可能成为 $i$ 的答案
        \item 更具体地，若 $\max(dis_i,dis_j)<\min_{i<k<j}dis_k$，则 $i,j$ 是有效点对
        \item 考虑如何找到这些点对
        \item 我们按照编号升序扫描分治块内的点，维护不降单调栈，弹栈时完成支配点对的配对。个数是 $O(n\log n)$
        \item 然后转化为一个二维排序的问题，类似于查询二维区域内最小值，随便解决即可
    \end{itemize}
\end{frame}

% ===== 结束页 =====

\begin{frame}
    \begin{center}
        {\Huge\calligra Thanks!}
    \end{center}
\end{frame}

\end{document}
